\documentclass{osucourse}
\course{1126}

%% Not in the university's course data, so these come
%% from the published sheet this file was imported from.
\SetCourseField{exclusions}{Not open to students with credit for 107.}

\begin{document}

\CatalogDescription
\PrerequisitesAndExclusions

\begin{Purpose}
To develop an appreciation of, and basic competency in, the use of analytical thought in the development of a cohesive body of useful mathematical knowledge, with special emphasis on topics encountered in elementary and middle school mathematics programs. Math 1126 addresses basic geometric concepts and measurement, symmetry and rigid motions, congruence, similarity and scaling, coordinate geometry, algebraic thinking, linear functions, counting techniques and probability. \emph{Appropriate for those preparing to become early childhood educators and for those preparing to teach subjects other than math in middle school}.
\end{Purpose}

\begin{Text}
\emph{Mathematics for Elementary Teachers, with Activity Manual}, 4\textsuperscript{rd} Edition, by Sybilla Beckmann, Pearson, ISBN for the package is 9780321836715 (loose-leaf) and Student Packet.
\end{Text}

\begin{TopicsList}
\item Spatial visualization and basic geometric concepts: angles, 2- and 3-D shapes and their properties.
\item Measurement: meaning of length, area, volume, measurement techniques, unit conversion, actions preserving area/volume, and scaling.
\item Transformations: symmetry, congruence, similarity.
\item Geometric constructions with various tools (compass, paper folding).
\item Algebraic thinking: expressions, measurement formulas, scaling, functions, use of formulas, graphs, and tables, sequences, and coordinate geometry.
\item Counting: inclusion/exclusion, fundamental counting principle, tree diagrams, permutations and combinations, Pascal’s triangle.
\item Basic ideas of probability: Law of Large Numbers, sample and event spaces, use of tree diagrams, simulations, and discussion of common misconceptions.
\item Problem solving and justifications at multiple levels are themes of the course. *Currently taught in either lecture/recitation or workshop format.
\end{TopicsList}

\end{document}
